\documentclass[answers]{exam}

% Version compilable
%\usepackage{../../../mypackages}
%\usepackage{../../../macros}

% Version pour execution du script
\usepackage{../../../../mypackages}
\usepackage{../../../../macros}

% Pour le tableau
\usepackage{array}
\usepackage{longtable}
\usepackage{multirow}
\usepackage[table]{xcolor}

\title{Notation des cahiers}
\author{Lutin + Cahier d'exercices de Salomé}
\date{2 Décembre 2025}

% Mise en forme des solutions en bleu
\SolutionEmphasis{\color{blue}}
\renewcommand{\solutiontitle}{\noindent}

\definecolor{lavande}{RGB}{180,190,255}

% Colonne p{} alignée à gauche et qui wrappe
\newcolumntype{L}[1]{>{\raggedright\arraybackslash}p{#1}}

\begin{document}

\textbf{Collège Lycée Suger}
\hfill
\textbf{Mathématiques} \\

\textbf{Année 2025-2026}
\hfill
\textbf{Classe de 6ème} \par

{\let\newpage\relax\maketitle}

\begin{tcolorbox}[
  %colback=blue!5!white,
  colframe=red!90!black,
  %title=Commentaire plus détaillé,
  fonttitle=\bfseries,
  boxsep=4pt,
  sharp corners,
]
\begin{center}
\textbf{\textcolor{red}{Ce document doit être imprimé et \\ (1) stocké dans le lutin OU (2) collé dans le cahier d'exercice.}}
\end{center}
\end{tcolorbox}

\subsection*{Note}

\renewcommand{\arraystretch}{1.7}
{
  %\setlength{\LTleft}{-2.3cm}  % déborde de 1.5 cm dans la marge de gauche
%\setlength{\LTright}{0pt}    % rien de plus à droite

\centering
\begin{longtable}{|L{0.1\textwidth}|L{0.5\textwidth}|c|>{\columncolor{green!35}}c|}
\hline
\rowcolor{lavande}
\textbf{Catégorie} & \textbf{Critère évalué} & \textbf{Points possibles} & \textbf{Points obtenus} \\[2pt]
\hline

% ---------------- Cahier d'exercices ----------------
\multirow{3}{*}{
  \parbox[t]{0.15\textwidth}{\raggedright Cahier d’exercices}}
  & Titres des exercices présents et soignés (+ utilisation de couleur)
  & 2
  & 1 \\
\cline{2-4}

  & Ordres de grandeurs notés
  & 1 
  & 1 \\
\cline{2-4}

  & Corrections notées (ou rattrapées en cas d'absence)
  & 3
  & 1.5 \\
\hline  

% ---------------- Lutin ----------------
\multirow{5}{*}{Lutin}
  & Feuilles rangées dans l’ordre
  & 1
  & 1 \\
\cline{2-4}

  & 1-2 feuilles max par pochette
  & 1
  & 1 \\
\cline{2-4}

  & Interrogations + calcul mental faciles à retrouver
  & 2
  & 1 \\
\cline{2-4}

  & Cours complet (rien de manquant)
  & 2
  & 1.5 \\
\cline{2-4}

  & Leçons, phrases, définitions recopiées correctement
  & 2
  & 0.75 \\
\hline

% ---------------- Général ----------------
\multirow{5}{*}{Général}
  & Organisation générale (titres soulignés, couleurs…)
  & 2
  & 1.5 \\
\cline{2-4}

  & Écriture soignée et lisible
  & 1
  & 0.25 \\
\cline{2-4}

  & Peu ou pas de ratures
  & 1
  & 0.25 \\
\cline{2-4}

  & Pas de feuilles volantes
  & 1
  & 1 \\
\cline{2-4}

  & Présence de prénom / nom // Identification facile du lutin
  & 1 
  & 1 \\
\hline

% ---------------- Total ----------------
\rowcolor{green!35}
 & \textbf{Total} & \textbf{20} & \textbf{ 12.75 } \\

\hline

\end{longtable}
}

\subsection*{Commentaire et conseils}

\begin{tcolorbox}[
  colback=blue!5!white,
  colframe=blue!70!black,
  title=Commentaire plus détaillé,
  fonttitle=\bfseries,
  boxsep=4pt,
  sharp corners,
]
\begin{compactitem}
  \item Fais bien ressortir les titres de chapitres et les énoncés des exercices dans ton cahier.
  \item Améliore la rédaction des corrections : écris des phrases complètes et claires.
  \item Termine les corrections commencées et veille à ce qu'elles soient toutes présentes.
  \item Continue à soigner ton travail, tu peux encore gagner en précision.
\end{compactitem}
\end{tcolorbox}



\subsection*{Rappel - Barème de notation des cahiers}

\renewcommand{\arraystretch}{1.7}

{
  \setlength{\LTleft}{-2.3cm}  % déborde de 1.5 cm dans la marge de gauche
\setlength{\LTright}{0pt}    % rien de plus à droite

\centering
\begin{longtable}{|L{0.1\textwidth}|L{0.3\textwidth}|L{0.7\textwidth}|c|}
\hline
\rowcolor{lavande}
\textbf{Catégorie} & \textbf{Critère évalué} & \textbf{Description} & \textbf{Points} \\[2pt]
\hline

% ---------------- Cahier d'exercices ----------------
\multirow{3}{*}{
  \parbox[t]{0.15\textwidth}{\raggedright Cahier d’exercices}
}
  & Titres des exercices présents et soignés (+ utilisation de couleur)
  & Les exercices comportent un titre clair, souligné ou encadré. Et une référence claire à l'endroit où retrouver l'exercice (« Exercice n°34 page 127 – Le livre scolaire », « Exercice d'application n°3 – Lutin », « Exercice n°2 – Fiche d'exercice »).
  & 2 \\
\cline{2-4}

  & Ordres de grandeurs notés
  & Les ordres de grandeurs donnés en début d'heure ont bien été notés.
  & 1 \\
\cline{2-4}

  & Corrections notées (ou rattrapées en cas d'absence)
  & Toutes les corrections sont recopiées, complétées ou rattrapées en cas d’absence, de manière lisible. Même si vous avez juste, j'attends de vous que vous preniez la correction. Pour progresser sur la rédaction.
  & 3 \\[6pt]
\hline  

% ---------------- Lutin ----------------
\multirow{5}{*}{Lutin}
  & Feuilles rangées dans l’ordre
  & Les chapitres sont dans l'ordre. Et les feuilles sont classées dans l’ordre attendu : Recto : Page 1. Verso : Page 2. Puis nouvelle pochette transparente : Page 3 au recto, Page 4 au verso etc.
  & 1 \\
\cline{2-4}

  & 1-2 feuilles max par pochette
  & Chaque feuille imprimée dans un feuillet transparent (1-2 feuilles par pochette transparente). Je ne veux pas voir toutes les feuilles d'un chapitre dans une seule et même pochette transparente.
  & 1 \\
\cline{2-4}

  & Interrogations + calcul mental faciles à retrouver
  & Les évaluations de calcul mental sont regroupées et rapidement retrouvables.
  & 2 \\
\cline{2-4}

  & Cours complet (rien de manquant)
  & Aucune feuille manquante : l’ensemble du cours est présent dans le lutin.
  & 2 \\
\cline{2-4}

  & Leçons, phrases, définitions recopiées correctement
  & Le travail est propre : tracés à la règle ou compas, titres alignés, soin général. La prise de notes est complète. Il ne manque pas de définition, de schéma.
  & 2 \\[6pt]
\hline

% ---------------- Général ----------------
\multirow{5}{*}{Général}
  & Organisation générale (titres soulignés, couleurs…)
  & Le cahier et le lutin sont organisés : titres visibles, couleur, structure claire.
  & 2 \\
\cline{2-4}

  & Écriture soignée et lisible
  & L’écriture est propre, lisible, régulière.
  & 1 \\
\cline{2-4}

  & Peu ou pas de ratures
  & Les ratures sont limitées et propres, sans nuire à la lecture.
  & 1 \\
\cline{2-4}

  & Pas de feuilles volantes
  & Aucune feuille glissée en vrac ; tout est rangé ou collé.
  & 1 \\
\cline{2-4}

  & Présence de prénom / nom // Identification facile du lutin
  & On sait très clairement à qui appartient le lutin, et à qui appartient le cahier d'exercice. Il y a une page de garde qui indique clairement qu'il s'agit de votre matériel de Mathématiques
  & 1 \\[6pt]
\hline

% ---------------- Total ----------------
\rowcolor{green!35}
 & \textbf{Total} & \textbf{Note maximale possible} & \textbf{20} \\
\hline

\end{longtable}
}



\subsection*{Conseils généraux}

\raggedright
\noindent\textbf{Gestion du cahier d'exercice}\par\vspace{0.5\baselineskip}

\begin{compactitem}
  \item A chaque début de chapitre, écrivez en gros, en rouge le titre du chapitre. Encadrez-le. Tout ce qui sera écrit en dessous seront des exercices en lien avec ce Chapitre
  \item Indiquez clairement à quel exercice une correction est associée. Les exercices peuvent venir de plusieurs endroits (Fiche d'exercice, Livre Scolaire, Polycopié de cours), donc mettre simplement "Exo 1" peut porter à confusion. On écrit, dans une couleur bien identifiable (rouge par exemple, souligné, ou encadré) : "Exercice N°26 Page 127 du Livre Scolaire", ou "Exercice d'application N°1 - Polycopié Lutin". Il faut que vous puissiez retrouver facilement l'exercice auquel vous faites référence.
  \item Vous pouvez même coller l'énoncé dans le cahier. 
  \item \textbf{\textcolor{red}{Même quand vous avez juste : si vous voyez que ma rédaction est plus complète, j'attends que vous la notiez. Il faut que vous progressiez en rédaction. Avoir la bonne réponse ne suffit pas.}}
\end{compactitem}

\begin{figure}[H]
  \centering
  \includegraphics[width=0.4\linewidth]{cahier_rempli.jpg}
  \captionsetup{labelformat=empty}
\end{figure}

\vspace{1em}
\noindent\textbf{Gestion du lutin}\par\vspace{0.5\baselineskip}

\begin{compactitem}
  \item Je ne veux pas voir des Chapitres dans le désordre, où le Chapitre 7 arrive avant le Chapitre 2.
  \item Respectez la numérotation des pages : Page 1 au recto, Page 2 au verso. Puis nouvelle pochette transparente : Page 3 au recto, Page 4 au verso etc. Il arrive souvent que les feuilles soient dans le désordre (Page 1, puis Page 6 au verso, puis Page 4 etc).
  \item L'imprimante a eu quelques problèmes, en imprimant dans l'autre sens des versos. Je vous conseille de ré imprimer les feuilles qui posent problème. Ou demandez-moi. 
  \item Je vous conseille de coller vos interrogations de calcul mental sur des feuilles simples ou doubles. Et stockez ces feuilles à la fin de votre lutin. Vous devez être capables de retrouver très rapidement toutes vos interrogations de calcul mental - les mettre au même endroit facilite les choses.
  \item Il est attendu que vos Devoirs sur table soient aussi stockés dans votre lutin (idéalement avec la correction détaillée imprimée) 
  \item Prise de note : 
  \begin{compactitem}
    \item Les figures doivent être faites au crayon à papier.
    \item Le reste de la prise de note se fait au stylo (je vois beaucoup de prise de note / théorèmes etc écrits au crayon à papier).
    \item Des légendes doivent figurer sur les figures.
    \item \textbf{On ne fait pas de tracé à main levée, on réalise ses figures avec compas, rapporteur, règle ou équerre. Forcez-vous à prendre de bonnes habitudes.}
  \end{compactitem}
\end{compactitem}

\vspace{1em}
\noindent\textbf{Général}\par\vspace{0.5\baselineskip}

\begin{compactitem}
  \item Evitez les ratures - vos cahiers sont parfois illisibles
  \item Soignez votre écriture
  \item Il ne faut pas qu'il y ait de feuille volante. Tout doit être soit collé, soit rangé dans le lutin. 
\end{compactitem}



\end{document}